\chapter{Introduction}




This part is devoted to studying curves and surfaces in $\mathbb{R}^3$. More precisely, this part is going to address the following questions:

\subsection*{Curves}

\begin{itemize}
\item \textbf{How to describe a curve?}

    e.g.: $\{ x = y, z = 0\}$ or equivalently $\{ \begin{pmatrix}
        t\\
        t\\
        0
    \end{pmatrix} \mid t \in \mathbb{R} \}$

\item \textbf{Properties of curves?}

    e.g.: length, curvature.

\item \textbf{If we know the curvature of a curve at every point, can we recover the whole curve?}

\item \textbf{If we suppose that we have a closed curve in $\mathbb{R}^2$ of a given length, what is the largest possible area bounded by the curve?} (Isoperimetric Inequality)
\end{itemize}

\subsection*{Surfaces}

\begin{itemize}
\item \textbf{How to describe a surface? More precisely, what kind of surfaces shall we study?}

    e.g.: We'd not learn surfaces with corners, self-intersection, and cusps (``infinite tangent'').
    
    \emph{Remark}: These kinds of surfaces will be discussed in the discrete Gauss--Bonnet Theorem.

\item \textbf{For the ``smooth'' surface, how can we study the ``area'' and ``curvature'' of them?}

\item \textbf{How to compute the ``shortest'' distance between two points of a surface?} (Geodesics)
\end{itemize}